\documentclass[10pt]{article}
\usepackage[letterpaper]{geometry}
\usepackage{hicss}
\usepackage{times}
\usepackage[none]{hyphenat}
\usepackage{url}
\usepackage{latexsym}
\usepackage{indentfirst}
\usepackage{graphicx}
\graphicspath{{images/}}
\usepackage{amsmath}
\usepackage{booktabs}
\usepackage{xcolor}
\usepackage{enumitem}
\setlist{noitemsep}
\usepackage{csquotes}
\usepackage[
    style=apa,
  ]{biblatex}
\addbibresource{references.bib}

\usepackage{amsmath,amssymb,amsfonts}
\usepackage{algorithmic}
\usepackage{algorithm}
\usepackage{graphicx}
\usepackage{textcomp}
\usepackage{xcolor}
\usepackage{hyperref}

\newcommand{\sansserifformat}[1]{\fontfamily{cmss}{ #1}}%

\setlength\titlebox{7cm}

\title{Agentic Swarm Coordination: Decentralized Target Allocation Using the OpenClaw Framework in D-DIL Environments}

\author{Aadi Sharma \\
 SRM Institute of Science and Technology \\
 {\underline{ as2382@srmist.edu.in}} \\ \And
 Hrishitva Patel \\
 Department of Information Systems and Cyber Security \\
 Carlos Alvarez College of Business, University of Texas at San Antonio \\
 {\underline{ hrishitva.patel@utsa.edu} } \\ }

\date{}

\begin{document}
\maketitle

\begin{abstract}
\textit{The deployment of autonomous drone swarms in Denied, Disrupted, Intermittent, and Limited (D-DIL) environments represents a formidable challenge in modern networked robotics. Under severe Electronic Warfare (EW) and Radio Frequency (RF) jamming, traditional centralized command-and-control architectures—which rely heavily on high-bandwidth cloud or edge servers—experience catastrophic failure due to extreme packet loss and network partitioning. To address this vulnerability, this paper proposes a novel system architecture for decentralized drone swarms acting as "frugal edge nodes." By integrating the open-source OpenClaw autonomous agent framework, our architecture bypasses the need for centralized intelligence. Instead, we propose a decentralized auction protocol modeled after consensus-based task allocation to achieve swarm-level optimization using a proxy heuristic function. This function dynamically evaluates target suitability based on local perception, battery life, kinematic distance, and payload capabilities. Communication efficiency is maximized through extreme data compression, broadcasting encrypted textual state representations rather than bandwidth-heavy video streams over a localized message bus. This paper presents the comprehensive theoretical framework, system architecture, and mathematical foundation of the proposed model. Furthermore, we outline an Object-Oriented Software Emulation methodology utilizing Python and the NS-3 network emulator to validate the architecture under simulated EW conditions.}
\end{abstract}

\subsubsection*{Keywords:}
Agentic AI, D-DIL Environments, Swarm Intelligence, Frugal Edge Computing, Decentralized Auction Protocol, OpenClaw, FANETs.

\section{Introduction}

The operational paradigm of uncrewed aerial systems (UAS) has increasingly shifted from single-agent, remotely piloted operations to multi-agent, autonomous swarm deployments. This evolution is predicated on the ability of flying ad-hoc networks (FANETs) to maintain robust, high-throughput communication links. However, in Denied, Disrupted, Intermittent, and Limited (D-DIL) environments—characteristically defined by adversarial Electronic Warfare (EW), severe Radio Frequency (RF) jamming, and physical obstructions—the foundational assumption of reliable communication architectures collapses.

\subsection{The Vulnerability of Centralized Paradigms}
Centralized artificial intelligence approaches often assume that high-fidelity telemetry, including multi-megapixel video feeds and complex sensor data, can be seamlessly streamed to a central server for processing. The server, equipped with significant computational resources, runs complex optimization algorithms and transmits actionable commands back to the drones. In D-DIL environments, where packet loss can routinely exceed 80\%, this centralized loop is completely untenable. The reliance on continuous, high-bandwidth telemetry results in an intolerable latency bottleneck. When the communication link to a central command node or a cloud processing server is severed, the swarm's collective intelligence degrades into disjointed, inefficient, and potentially catastrophic local behaviors.

\subsection{Introducing the Frugal Edge}
To circumvent this architectural bottleneck, we propose a paradigm shift toward the "Frugal Edge." The Frugal Edge leverages aggressive localization of computational reasoning and perception, treating each drone as a sovereign, intelligent node. By stripping away the dependency on large-scale centralized models, we empower the swarm with Agentic AI—specifically leveraging the OpenClaw framework—to conduct high-level reasoning onboard resource-constrained hardware. By migrating the intelligence from the cloud to the drone, we drastically mitigate the need for constant, high-bandwidth communication. 

\subsection{Contributions of this Paper}
The core contribution of this paper is the formalization of a decentralized system architecture that enables target allocation without a central server. This is achieved through a localized, consensus-based auction protocol where individual drones calculate bidding metrics using a proxy heuristic function based on their immediate state variables. This paper details the structural foundations of this architecture, establishing the mathematical formulations required to evaluate its efficacy without relying on fragile external network dependencies.

\section{Literature Review}

The conceptual evolution of distributed intelligence at the network periphery has been extensively documented, transitioning from passive data collection architectures to highly active, autonomous frameworks capable of operating in hostile environments.

\subsection{Edge Computing and Military IoT}
The genesis of edge computing emerged as a solution to the latency and bandwidth constraints inherent in cloud-centric architectures. Shi et al. \cite{shi2016} established the foundational vision for edge computing, defining it as the processing of data at the periphery of the network, close to the data sources. This paradigm shift was primarily aimed at commercial Internet of Things (IoT) applications, where passive sensors generated massive volumes of data that required local filtering before cloud transmission. 

However, the battlefield environment imposes significantly more stringent constraints than commercial settings. Suri et al. \cite{suri2016} analyzed the applicability of IoT architectures to military operations, coining the concept of the Internet of Battlefield Things (IoBT). They highlighted that unlike commercial IoT, military networks operate under adversarial conditions, where nodes are subject to kinetic destruction, and communication channels face active EW disruption.

\subsection{Flying Ad-Hoc Networks (FANETs) Constraints}
The transition from static ground sensors to mobile platforms necessitated the development of Flying Ad-Hoc Networks (FANETs). Bekmezci et al. \cite{bekmezci2013} provided a comprehensive survey of FANETs, distinguishing them from traditional Mobile Ad-Hoc Networks (MANETs) due to their higher node mobility, 3D spatial distribution, and unique topological volatility. FANETs are highly susceptible to network partitioning, rendering traditional routing protocols obsolete when scaling beyond tightly clustered formations.

\subsection{Federated Learning at the Extreme Edge}
As the demand for decentralized intelligence grew, Federated Learning (FL) emerged as a mechanism to train machine learning models across distributed nodes without sharing raw data. McMahan et al. \cite{mcmahan2017} introduced the foundational algorithms for communication-efficient learning from decentralized data. Kairouz et al. \cite{kairouz2021} later expanded on the open problems within FL, specifically pointing to the challenges of statistical heterogeneity and communication bottlenecks. 

Recent advancements have focused on resource-constrained IoT devices. Imteaj et al. \cite{imteaj2022} surveyed FL techniques optimized for microcontrollers and low-power devices. Building upon this, Junior et al. \cite{junior2025} proposed the FedSensor framework, pushing FL to the extreme edge to ensure secure, privacy-preserving sensor inference. Despite these advancements, FL remains primarily a model-training paradigm rather than a real-time, zero-shot operational reasoning framework required for dynamic target allocation.

\subsection{Agentic AI and Consensus-Based Auctions}
The shift from distributed model training to distributed autonomous action requires active agentic frameworks capable of reasoning, planning, and executing sequences of actions. The OpenClaw framework, developed by Steinberger \cite{steinberger2026}, provides a lightweight, open-source architecture for deploying personal AI assistants and autonomous agents. OpenClaw allows for the execution of complex logical chains and tool utilization on local hardware.

For swarm coordination, reasoning must be coupled with an allocation mechanism. Choi et al. \cite{choi2009} introduced the Consensus-Based Bundle Algorithm (CBBA), a robust, decentralized auction protocol for task allocation. CBBA guarantees conflict-free assignments by utilizing local bidding based on marginal scores and a consensus mechanism over a dynamic communication topology. Our proposed architecture deeply integrates the philosophical underpinnings of CBBA with the reasoning capabilities of OpenClaw.

\section{System Architecture: Perception and Reasoning}

The proposed architecture is specifically designed to eliminate the single point of failure inherent in centralized swarm command. By treating each drone as an autonomous, self-contained Frugal Edge Node, the system can dynamically optimize target allocation through a localized message bus, even when communication topology is severely fragmented.

\subsection{The Local Perception Node: Extreme Data Compression}
In conventional architectures, drones act primarily as sensor platforms, streaming high-resolution video to a central unit. A 1080p video stream compressed via H.264 typically requires between 3 to 6 Mbps. In a D-DIL environment, where the total available swarm bandwidth might drop below 50 Kbps, streaming raw sensor data is mathematically impossible.

Our architecture mandates that all sensor processing occurs strictly locally on the physical drone hardware. The Local Perception Node utilizes highly optimized, quantized convolutional neural networks to parse incoming sensor data in real-time. Instead of transmitting imagery, the perception node extracts semantic meaning and compiles it into an extremely compact, structured text string.

For example, a drone identifying a target converts the visual data into a serialized vector containing Target ID, Type classification, GPS coordinates, estimated velocity, and model confidence score. This payload is reduced to approximately 40 bytes. This extreme data compression enables the swarm to share target awareness using highly robust, low-bandwidth modulation schemes.

\subsection{The OpenClaw Reasoning Engine: Proxy Heuristic Function}
Once a target state vector is acquired or received from a peer, it is passed to the local OpenClaw Reasoning Engine \cite{steinberger2026}. In this architecture, OpenClaw operates as a deterministic, mathematically grounded decision engine to calculate a Suitability Score ($S_i(t_j)$) for drone $i$ relative to target $t_j$. 

To bypass the computational complexity of solving NP-hard global optimization problems onboard a microcomputer, OpenClaw evaluates a proxy heuristic function. This function locally evaluates the viability of an engagement based on three critical intrinsic variables: Battery Life ($B_i$), Kinematic Distance ($D_{ij}$), and Payload Match ($P_{ij}$).

The Suitability Score is formulated as:
\begin{equation}
S_i(t_j) = \max \left(0, \alpha \hat{B}_i - \beta \hat{D}_{ij} + \gamma P_{ij} \right)
\label{eq:suitability}
\end{equation}

Where $\hat{B}_i$ is the normalized remaining battery life, $\hat{D}_{ij}$ is the path-planning cost, $P_{ij}$ is the payload matching coefficient, and $\alpha, \beta, \gamma$ are tunable hyperparameters dynamically adjusted by the OpenClaw engine.

